For a future cash dividend, the model needs the bank's declaration as well as
the eventual payment. The OCC's national-bank dividend rules, retrieved for
1~January 2026, use the declaration date to determine compliance. They define
retained net income as net income less all dividends declared in the period.
Absent OCC approval under the earnings provision, current-year declarations, including
the proposal, are compared with current-year net income to date plus the
previous two years' retained net income, less specified required transfers.
The current year is the calendar year of declaration. A cash-flow statement
of dividends paid therefore cannot supply the declaration history by itself
\citep[12~CFR~5.61--5.64]{occnationaldividends2026}.

The preceding two years are not always enough. Section~5.64(c)(2) provides
conditional offsets for retained losses caused by dividends exceeding a
single year's net income. Those offsets can require the third and fourth
preceding years and follow a specified order; other earnings deficits do not
qualify for that adjustment. An implementer must preserve the cause of the
loss and the allocation of offsets before computing the earnings restriction.
Capital, undivided-profit and supervisory conditions remain separate: an
earnings calculation alone cannot establish a permissible transfer
\citep[12~CFR~5.63--5.65]{occnationaldividends2026}.

Selecting the applicable capital rule also requires more than a calendar
lookup. The final leverage rule published on 1~December 2025 is effective
1~April 2026 but permits early adoption from 1~January
\citep[p.~55248]{eslrfinal2025}. JPMorgan's 2025
Form~10-K reports that the Firm elected early adoption effective 1~January
2026 \citep[p.~91]{jpmorgan2025tenk}.
The rule's publication, general effective date, permitted election date and
the issuer's reported election answer different questions. The later filing
provides evidence about the election; it does not establish that the election
was publicly known on 1~January. A historical information query must retain
that distinction even when reconstructing an earlier regulatory state.

The amendment removes the additional 6\% supplementary-leverage condition
from the affected national banks' well-capitalized classification and adds
separate leverage-buffer payout provisions. That classification change does
not remove the distinct undercapitalization test or establish unrestricted
distributions \citep[pp.~55287--55288]{eslrfinal2025}. The reported Firm election
also does not supply every bank-specific application fact. A permissible
dividend calculation still needs the relevant capital definitions, distribution
limits, supervisory conditions and bank-level inputs. The US\$30.1bn remains
the issuer's dated capacity disclosure, rather than a recomputed legal limit.

